\section{Edge Pipeline Architecture}
The X-IDS pipeline is designed to simulate a real-world deployment where network traffic is captured at a gateway and processed locally.

\subsection{Distributed Data Stream (Kafka)}
We utilize **Apache Kafka** as a low-latency message broker. Network flow summaries (JSON) are produced at the capturing node and consumed by the Raspberry Pi. Kafka ensures that even if the inference engine experiences temporary spikes in latency, the data stream is buffered and not lost.

\subsection{Distributed Inference (PySpark on RPi)}
The core of the system is a **PySpark Streaming` application running on the Raspberry Pi 4B. 
\begin{itemize}
    \item \textbf{Micro-Batching}: Spark processes flows in small batches (e.g., 5-10 flows) to balance throughput and latency.
    \item \textbf{Performance Optimization}: We disable the Spark Web UI and limit the number of worker cores to two, ensuring that sufficient background resources are available for the OS and Kafka consumer.
\end{itemize}

\subsection{Storage and Alerting Stack}
Detected attacks are routed to two separate systems:
\begin{enumerate}
    \item \textbf{InfluxDB}: A time-series database for high-velocity metrics (throughput, latency, attack counts) visualized in Grafana.
    \item \textbf{PostgreSQL}: Stores persistent alert logs, including the SHAP explanation for each detected attack.
\end{enumerate}
